\documentclass[11pt,a4paper]{article}

% -------------------------------------------------
% Packages
% -------------------------------------------------
\usepackage[utf8]{inputenc}
\usepackage[T1]{fontenc}
\usepackage[brazil]{babel}
\usepackage[a4paper,top=1.2cm,bottom=1.2cm,left=1.4cm,right=1.4cm]{geometry}
\usepackage{hyperref}
\usepackage{xcolor}
\usepackage{titlesec}
\usepackage{enumitem}
\usepackage{tabularx}
\usepackage{fontawesome5}

\hypersetup{
    colorlinks=true,
    urlcolor=black
}

\definecolor{primary}{RGB}{25,55,109}

\pagestyle{empty}
\setlength{\parindent}{0pt}

\titleformat{\section}
{\Large\bfseries\color{primary}}
{}
{0pt}
{}
[\titlerule]

% Espaço antes/depois dos títulos de seção.
% Se passar de 2 páginas, diminua (ex.: {10pt}{5pt}); se sobrar espaço, aumente.
\titlespacing*{\section}{0pt}{14pt}{7pt}

\begin{document}

%%%%%%%%%%%%%%%%%%%%%%%%%%%%%%%%%%%%%%%%%%%%%%%%%%
% HEADER
%%%%%%%%%%%%%%%%%%%%%%%%%%%%%%%%%%%%%%%%%%%%%%%%%%

\begin{center}

{\Huge\bfseries Rafael Caetano da Silva}

\vspace{5pt}

{\Large Desenvolvedor .NET Pleno}

\vspace{4pt}

{\small
C\# • .NET • ASP.NET Core Web API • Entity Framework Core • SQL Server • PostgreSQL • Azure DevOps
}

\vspace{6pt}

\small

\faEnvelope\
\href{mailto:rafaelkaetanodf@gmail.com}{rafaelkaetanodf@gmail.com}
\quad|\quad
\faPhone\
(61) 99398-9820
\quad|\quad
\faMapMarker*
Brasília - DF

\vspace{3pt}

\faGithub\
\href{https://github.com/Rafael-CaetanoS}{github.com/Rafael-CaetanoS}
\quad|\quad
\faLinkedin\
\href{https://linkedin.com/in/rafaelcaetanoo}{linkedin.com/in/rafaelcaetanoo}

\end{center}

\vspace{-0.2cm}
%%%%%%%%%%%%%%%%%%%%%%%%%%%%%%%%%%%%%%%%%%%%%%%%%%
\section{Resumo Profissional}

Desenvolvedor .NET com experiência em C\#, .NET e ASP.NET Core Web API no desenvolvimento de integrações entre sistemas e aplicações corporativas. Atuação na construção e consumo de APIs REST, integração com bancos de dados relacionais (SQL Server, PostgreSQL, MySQL e Oracle) utilizando SQL e Entity Framework Core, mensageria com filas no Azure Service Bus para processamento assíncrono e desacoplamento entre serviços, e comunicação em tempo real via WebSockets com SignalR.

\vspace{0.15cm}

Participo da definição da arquitetura das soluções, do levantamento de requisitos técnicos junto às áreas de negócio, da revisão de código via pull requests e da tomada de decisões técnicas, exercendo atualmente funções de liderança técnica. Experiência com Clean Architecture, Domain-Driven Design (DDD), CQRS, arquitetura em camadas, SOLID, injeção de dependência e Design Patterns, além de testes automatizados com xUnit e pipelines de CI/CD no Azure DevOps.

%%%%%%%%%%%%%%%%%%%%%%%%%%%%%%%%%%%%%%%%%%%%%%%%%%
\section{Experiência Profissional}

\textbf{Coordenador Back-end}
\hfill Ago/2025 -- Atual

Grupo Orion

\begin{itemize}[leftmargin=*,noitemsep,topsep=2pt]

\item Desenvolvimento e manutenção de integrações entre sistemas corporativos, serviços externos e plataformas de monitoramento IoT utilizando C\#/.NET e ASP.NET Core Web API.

\item Construção e evolução de APIs REST utilizando ASP.NET Core Web API.

\item Revisão de código e de pull requests no Azure DevOps, validação de soluções técnicas e apoio aos desenvolvedores na resolução de problemas.

\item Definição da arquitetura de novas funcionalidades utilizando Clean Architecture, Domain-Driven Design (DDD) e princípios SOLID.

\item Tomada de decisões técnicas relacionadas à arquitetura, banco de dados, integrações, desempenho e escalabilidade das aplicações.

\item Liderança técnica da equipe Back-end, apoiando o desenvolvimento dos profissionais e promovendo boas práticas de engenharia de software.

\item Levantamento e análise de requisitos técnicos junto às áreas de negócio e produto, planejamento técnico das demandas e distribuição das atividades da equipe.

\item Estruturação e manutenção de pipelines de CI/CD (integração e entrega contínua) utilizando Azure DevOps.

\end{itemize}

\vspace{0.2cm}

\textbf{Desenvolvedor Back-end Pleno}
\hfill Abr/2025 -- Ago/2025

Grupo Orion

\begin{itemize}[leftmargin=*,noitemsep,topsep=2pt]

\item Desenvolvimento de APIs REST e integrações entre sistemas com C\# e ASP.NET Core Web API, seguindo Clean Architecture e DDD.

\item Modelagem e manutenção de bancos de dados relacionais PostgreSQL e MySQL.

\item Implementação de testes automatizados (testes unitários) utilizando xUnit para aumento da confiabilidade das aplicações.

\item Aplicação de SOLID, injeção de dependência (Dependency Injection) e Design Patterns para desenvolvimento de soluções de fácil manutenção.

\end{itemize}

\vspace{0.2cm}

\textbf{Analista de Sistemas (Estágio)}
\hfill Dez/2023 -- Dez/2024

Infraero

\begin{itemize}[leftmargin=*,noitemsep,topsep=2pt]

\item Desenvolvimento e manutenção de aplicações corporativas e sistemas internos de operações aeroportuárias com C\#, ASP.NET Core, Blazor e Java.

\item Desenvolvimento com bancos de dados relacionais SQL Server, Oracle, PostgreSQL e MySQL.

\item Versionamento de código com Git e Azure DevOps; Docker para padronização dos ambientes de desenvolvimento.

\item Desenvolvimento de relatórios utilizando Microsoft Report Builder.

\end{itemize}

%%%%%%%%%%%%%%%%%%%%%%%%%%%%%%%%%%%%%%%%%%%%%%%%%%
\section{Projetos Relevantes}

\textbf{Integração Infraspeak}

Desenvolvimento de integrações entre sistemas via APIs REST para sincronização de ativos, ordens de serviço, falhas e propriedades adicionais entre sistemas corporativos.

\vspace{0.2cm}

\textbf{Plataforma Cortex}

Responsável pelo desenvolvimento do backend da plataforma de monitoramento IoT, incluindo gerenciamento de alertas, notificações, configurações de alarmes e integração com dispositivos inteligentes.

\vspace{0.2cm}

\textbf{SIVWEB} -- Sistema de gerenciamento de monitores aeroportuários desenvolvido com ASP.NET Core e Blazor.

\vspace{0.2cm}

\textbf{SICOA} -- Sistema corporativo para identificação e controle de acesso de usuários.

%%%%%%%%%%%%%%%%%%%%%%%%%%%%%%%%%%%%%%%%%%%%%%%%%%
\section{Competências Técnicas}

\textbf{Back-end:} C\#, .NET, ASP.NET Core Web API, APIs REST, HTTP, Programação Assíncrona (async/await), Java, Spring Boot

\vspace{0.12cm}

\textbf{Banco de Dados:} SQL, SQL Server, PostgreSQL, MySQL, Oracle

\vspace{0.12cm}

\textbf{ORM:} Entity Framework Core, Dapper, Hibernate

\vspace{0.12cm}

\textbf{Mensageria e Tempo Real:} Mensageria, Filas, Azure Service Bus, SignalR, WebSockets

\vspace{0.12cm}

\textbf{Arquitetura:} Clean Architecture, Domain-Driven Design (DDD), CQRS, Arquitetura em Camadas, SOLID, Injeção de Dependência, MVC

\vspace{0.12cm}

\textbf{Testes:} Testes Automatizados, Testes Unitários, xUnit

\vspace{0.12cm}

\textbf{Cloud e DevOps:} Azure DevOps, Git, Pull Requests, Code Review, CI/CD (Integração e Entrega Contínua), Azure, Docker

\vspace{0.12cm}

\textbf{Front-end:} Angular, TypeScript, JavaScript, HTML, CSS

\vspace{0.12cm}

\textbf{Ferramentas:} OpenAPI/Swagger, Postman, Report Builder

%%%%%%%%%%%%%%%%%%%%%%%%%%%%%%%%%%%%%%%%%%%%%%%%%%
\section{Formação}

\textbf{Pós-graduação em Engenharia de Software}
\hfill 2025 -- 2026

Universidade Católica de Brasília

\vspace{0.2cm}

\textbf{Bacharelado em Sistemas de Informação}
\hfill 2021 -- 2024

Centro Universitário do Distrito Federal (UDF)

%%%%%%%%%%%%%%%%%%%%%%%%%%%%%%%%%%%%%%%%%%%%%%%%%%
\section{Idiomas}

Português — Nativo \quad|\quad Espanhol — Intermediário \quad|\quad Inglês — Básico

\end{document}
